% Terminal apparatus of the concise study: the scope note and the References
% for the sources this companion actually uses, at the loci it cites. The
% revision timestamp and the rights colophon follow from the shared
% generation record.

\runninghead{Scope and References}
\section{Appendix: Scope and Qualifications}

{\footnotesize
\textbf{Edition, Mass and occurrence.} The Mass treated is \latin{In festo
omnium Sanctorum}, \latin{Die 1 novembris}, I classis, of the \work{Missale
Romanum}, editio typica, Typis Polyglottis Vaticanis 1962, printed
pp.~702--703, nos.~3974--3983, as said on Sunday 1~November 2026 in the
universal 1962 calendar, with the commemoration of the Twenty-third Sunday
after Pentecost, pp.~416--417, nos.~1720, 1726 and 1728. The precedence of
the feast, the commemoration and its order rest on the Missal's general
rubrics, as \texttt{research/context.md} records. The feast has no Preface of
its own; the general rubrics (RGMR 482, 494~b, 498) admit both the Preface of
the Most Holy Trinity and the common Preface, and neither text is given here.
The Ordinary and the Office of the feast are not treated.

\textbf{Text and English.} The Latin was read on page images of the CMAA
facsimile of the 1962 typical edition and collated with the Tours printing of
1922 and the Ratisbon printing of 1862; the collation is in
\texttt{propers/verified.md}. The scriptural English is the Douay--Rheims in
Challoner's revision, and psalms are cited in the Missal's numbering. The
English of the Introit antiphon and of the Offertory is the Lasance
\work{New Roman Missal} of 1945, quoted on its own attribution. None of this
English is an approved liturgical translation of the 1962 Missal. Patristic
and later authors are quoted in the nineteenth-century English of the
editions named below, or in their Latin or French; where an edition has no
English, the English beside the quotation is the study's own rendering, and
no English is composed for a liturgical text.

\textbf{Reception.} The three interpretations, their sources and their four
senses are developed in full in the expansive study of this same feast, which
this companion condenses; its loci, editions and bounds are recorded in
\texttt{research/scope.md} and \texttt{research/interpretations.md}. The
reading is bounded by the corpora named there. Not read: Cassiodorus,
Theodoret and St Athanasius on the psalms; St Jerome on Mt~5; St Thomas
Aquinas on Mt~11; St Gregory of Nyssa on the Beatitudes; and the
\work{Glossa ordinaria}, known only as St Bonaventure reports it. Each
interpretation's principal authorities are Fathers and saints with a row in
the repository's register of authorial standing. St Bernard, St Bonaventure,
St Alphonsus Liguori, St Victorinus of Pettau, St Cyprian and St Bruno of
Würzburg have no row there and are cited in support only. Blessed Ildefonso
Schuster, Lucien Fromage and William Durandus are read on this Mass as a
Mass; Fromage and Durandus are ecclesiastical writers, and none of the three
is a principal authority for an interpretation. St Rabanus Maurus is a
compiler, credited only with the words his commentary marks as his own.
Every connection drawn between texts of the Mass belongs to the
interpretations unless an author is named as making it, and the orations of
the commemorated Sunday are read only for what they say where they stand in
this Mass. Augustine's \work{Enarrationes} are read in Migne's Latin; St
Robert Bellarmine in O'Sullivan's abridged English, which numbers Psalm~33
one verse lower than the Vulgate; St Basil in the Latin version printed in
PG~29; Aquinas on the Psalms and Bonaventure on uncorrected optical text,
checked at the loci marked in \texttt{research/scope.md}.

\textbf{Chronology.} Every biblical date on page~2 is generated from the
repository's Scripture chronology corpus, whose sources are the
\work{Catholic Encyclopedia} articles named below, Haydock's Douay commentary
and the introduction to the Psalms of the \work{New American Bible, Revised
Edition}. The record holds no date for the Sermon on the Mount or for the
occasion of Psalm~32. Each passage's fuller evidence is given in the expansive
study's Scriptural Date and Location.

\textbf{Rights and use.} The Latin is published on 17 U.S.C. 103(b) together
with its public-domain witnesses, the Mame printing of 1922 and the Pustet
printing of 1862, on the pre-1931 line recorded in \texttt{guidance/missals.md},
applied as of 26~October 2026. The Douay--Rheims and the Lasance English
(public domain in the United States by nonrenewal) are quoted from registered
editions. The other quotations are short and come from editions published
before 1931. Scripture, liturgical texts and quoted translations are not
relicensed by the repository's licence. This is a study aid to the 1962
Missal, not an official liturgical book, translation or critical edition;
this companion's production record is \texttt{research/production-review.md}.
\par}

\section{References}

{\footnotesize
\begin{itemize}[leftmargin=1.2em,itemsep=.15em]
\item \work{Missale Romanum ex decreto Sacrosancti Concilii Tridentini
restitutum, Summorum Pontificum cura recognitum}, editio typica (Typis
Polyglottis Vaticanis, 1962), CMAA facsimile: \work{Rubricae generales} and
\work{Rubricae generales Missalis romani}; Twenty-third Sunday after
Pentecost, pp.~416--417, nos.~1720, 1726, 1728; \work{Praefationes},
p.~293, no.~1082, and p.~297, no.~1086; Sanctorale, pp.~702--703,
nos.~3974--3983.
\item \work{Missale Romanum}, editio IV iuxta typicam (Tours: Mame, 1922),
pp.~748--750; \work{Missale Romanum} (Ratisbon: Pustet, 1862),
pp.~357--358.
\item Lasance, F.~X., \work{The New Roman Missal in Latin and English},
revised edition (New York: Benziger Brothers, 1945), pp.~1265--1268 (All
Saints).
\item \work{The Holy Bible Translated from the Latin Vulgate} (Douay--Rheims,
Challoner revision), Project Gutenberg text: 1~Kings 21:10--15; Ps 32:1;
33:1, 10--11; 90:1; 118:120; Wis 2:20; 3:1--3; Isa 11:2--3; 33:6; Ezech 9:4;
Mt 4:23--25; 5:1--12a; 11:28--29; Lk 13:23--24; Jn 10:28; Rom 9:4--5;
1~Cor 4:11; 13:12; 2~Cor 1:7; Gal 3:29; Phil 4:11--12; Apoc 1:9; 7:2--12.
\item Alphonsus Liguori, \work{Preparation for Death}, Ninth Consideration
(Dublin: Duffy, 1843), pp.~117--122.
\item Augustine, \work{Enarrationes in Psalmos} 32 (enarration 1), 33
(sermo~2, nn.~14--15) and 87 (n.~5), Migne Latin text (Wikisource).
\item Augustine, \work{De sermone Domini in monte} I.3.10--I.4.12; Sermons
XIX [LXIX Ben.] and LXI, NPNF, first series, vol.~6; \work{Enchiridion} 29,
NPNF, first series, vol.~3.
\item Basil the Great, \work{Homilia in Psalmum XXXII}, n.~1, and
\work{Homilia in Psalmum XXXIII}, nn.~6--7, PG 29 (Paris, 1860), Latin
version.
\item Bede, \work{Explanatio Apocalypsis}, book I, on 7:1--12, in E.~Marshall's
English, \work{The Explanation of the Apocalypse} (Oxford and London: James
Parker, 1878), pp.~44--46, 52--53.
\item Bellarmine, Robert, \work{Explanatio in Psalmos}, in John O'Sullivan's
English, \work{A Commentary on the Book of Psalms}, on Ps 33.
\item Bernard of Clairvaux, \work{In festo omnium Sanctorum} sermones 1, 2
and 5, PL 183 (Paris, 1854), cols.~453--478.
\item Bonaventure, \work{Commentarius in librum Sapientiae}, c.~3, in
\work{Opera omnia}, vol.~6 (Quaracchi, 1893), pp.~125--126.
\item Bruno of Würzburg, \work{Expositio Psalmorum}, on Ps 32:1, PL 142.
\item Chrysostom, John, \work{Homilies on Matthew} 15 and 38, NPNF, first
series, vol.~10.
\item Cyprian, \work{Ad Fortunatum} (\work{Exhortation to Martyrdom}) 11,
ANF, vol.~5.
\item Gregory the Great, \work{Homiliae in Evangelia} 34, n.~11 (Latin,
Wikisource).
\item Hilary of Poitiers, \work{In Evangelium Matthaei commentarius}, cap.~IV,
nn.~1--7, PL 9 (Paris, 1844), cols.~933--935.
\item Rabanus Maurus, \work{Expositio in Matthaeum} II, on 5:1--12, PL 107
(Paris, 1864), cols.~794--798.
\item Thomas Aquinas, \work{Super Evangelium S.~Matthaei lectura}, cap.~5, in
\work{Opera omnia}, vol.~3 (Venice: Bettinelli, 1745), pp.~67, 72;
\work{In Psalmos Davidis expositio}, on Pss 32 and 33, in \work{Opera omnia},
vol.~14 (Parma, 1863); \work{Summa theologiae} I q.~23 a.~7, in the English
Dominican translation.
\item Victorinus of Pettau, \work{Commentary on the Apocalypse}, on 7:2--9 and
12:6, ANF, vol.~7.
\item Anonymous, \work{Expositio in Apocalypsim}, hom.~6, in the appendix to
Augustine, PL 35 (Paris, 1845), cols.~2425--2426.
\item Durandus, William, \work{Rationale divinorum officiorum} VII.34, in the
1568 printing with Beleth, p.~799.
\item Fromage, Lucien, continuing Prosper Guéranger, \work{L'Année
liturgique}, vol.~15 (Paris: Oudin, 1909), \latin{La Toussaint}, \latin{A la
Messe}.
\item Schuster, Ildefonso, \work{The Sacramentary (Liber Sacramentorum)},
vol.~V (London: Burns Oates \& Washbourne, 1930), pp.~210--211.
\item \work{The Catholic Encyclopedia} (New York: Robert Appleton, 1907--1912),
for the page-2 chronology: Corbett, ``David'' (vol.~4); Durand, ``The New
Testament'' (vol.~14); Gigot, ``Book of Wisdom'' (vol.~15); Jacquier,
``Gospel of St.~Matthew'' (vol.~10); Van den Biesen, ``Apocalypse'' (vol.~1).
\item Haydock, George Leo, \work{The Holy Bible} with Haydock's commentary
(Loreto reprint, 2014), the chronological note at 1~Kings 21 reporting
Ussher.
\item \work{New American Bible, Revised Edition}, introduction to the Psalms
(USCCB website), summarized.
\end{itemize}
\par}
